\documentclass[11pt,a4paper]{article}
\usepackage[T1]{fontenc}
\usepackage[utf8]{inputenc}
\usepackage{newtxtext,newtxmath} 
\usepackage[mode=text]{siunitx}  
\usepackage[hmargin=2.2cm,vmargin=1.8cm]{geometry}

\pagestyle{empty}
\begin{document}
	
\noindent{\large\bfseries Six-month progress report, June--December 2026}\\[1ex]
Rasika Dilhani, PhD candidate, Victoria University of Wellington\\
Supervisor: Alejandro C. Frery
	
\subsection*{Overview}
My thesis, ``Ordinal Pattern Features: Statistical Properties and Their Applications to Time Series Clustering'', builds on the method of Bandt and Pompe (2002). From the distribution of ordinal patterns in a time series I compute information-theory quantifiers, namely the Shannon, Rényi and Tsallis entropies, the Fisher information and the statistical complexity, and I use them as features to characterise and cluster time series.
In this period I worked on two applications, Sentinel-1 SAR data and seismic amplitude at Whakaari/White Island, and on the simulation study and the writing of the thesis.
	
\subsection*{SAR land cover classification}
With co-authors from IIT Indore, India, and Universidad Nacional de Hurlingham, Argentina, including my secondary supervisor, Prof. Andrea Rey, I wrote ``Land Cover Classification in SAR Time Series Using Ordinal Pattern Features'', accepted at IEEE MIGARSS-InGARSS 2026. We used Sentinel-1 VH time series (2017--2023) from five land cover classes near Haldwani, India, combined each $2\times2$ pixel neighbourhood into one series, and computed four features with $D=5$: permutation entropy, statistical complexity, Fisher information and Fisher complexity.
An SVM classified Forest and Agriculture best (sensitivities of \SI{92}{\percent} and \SI{83}{\percent}), while Urban and Water were hard to separate.
I wrote the R code that builds the series and computes the features, prepared the figures and typeset the manuscript.

	
\subsection*{Seismic amplitude at Whakaari/White Island}
I retrieved 10-minute RSAM data from stations WIZ and WSRZ for July 2014 to December 2019 and wrote R code for preprocessing, feature computation and plotting.
	The features are computed over 7-day windows, moved forward either one day (6-day overlap) or six days (1-day overlap).
I clustered the weekly features with $k$-means and DBSCAN, trained a class-weighted support vector machine with an RBF kernel, and compared the results with the GeoNet Volcanic Alert Level record using PCA and Sankey diagrams.
The clusters did not separate the eruptive periods, including the eruption of 9~December 2019, although these periods stand out in the entropy--complexity plane and in the feature time series.
I am therefore testing a threshold rule that flags anomalous weeks, followed by event coincidence analysis (ECA) to check whether flagged weeks coincide with eruption weeks more often than expected by chance.
I summarised the clustering results in a report, ``Ordinal-pattern cluster analysis of RSAM time series for detecting volcanic seismic anomalies at Whakaari, 2014--2019'', and sent it to a volcanologist, who is reviewing the labels we use as ground truth.
	
\subsection*{Thesis}
The literature review is mostly written.
It shows that ordinal-pattern features are used in many fields but rarely for clustering time series or for event coincidence analysis.
In the simulation study I generate time series as convex combinations of ARMA, AR, MA, logistic-map and sine processes, with several mixing weights and sample sizes, compute 13 ordinal-pattern features, and check whether the entropy--complexity plane, $k$-means and DBSCAN group the series by their generating model.
I am now writing the methodology chapter; at present it is illustrated only with the bearing-fault data from my proposal, and I am adding the completed work to it.
	
\subsection*{Next six months}
I will finish the threshold and ECA analysis for Whakaari and submit it to a journal, finish the Wayanad landslide study, and add the simulation results to the methodology chapter.
I aim to have a complete draft of the thesis by April 2027.
Scholarship support would let me finish this work and submit on time.
	
\end{document}